% Part: first-order-logic
% Chapter: natural-deduction
% Section: soundness

\documentclass[../../../include/open-logic-section]{subfiles}

\begin{document}

\iftag{FOL}
      {\olfileid[hi]{fol}{ntd}{sou}}
      {\olfileid[hi]{pl}{ntd}{sou}}
      
\olsection{सुदृढ़ता}

\begin{explain}
स्वाभाविक निगमन जैसा !!^a{derivation} तंत्र \emph{सुदृढ़} है यदि वह ऐसी बातें
!!{derive} नहीं कर सकता जो वास्तव में तार्किकतः निष्पन्न नहीं होतीं। इस प्रकार
सुदृढ़ता, !!{derivation} तंत्रों के लिए एक प्रकार का सुनिश्चित सुरक्षा-गुण है।
जिस प्रमाण-सैद्धांतिक गुण पर विचार हो, उसके अनुसार हम उदाहरणार्थ यह जानना
चाहेंगे कि
\begin{enumerate}
\item हर !!{derivable} !!{sentence} \iftag{FOL}{वैध}{सर्वसत्य} है;
\item यदि !!a{sentence} कुछ अन्य वाक्यों से !!{derivable} है, तो वह उनका
  तार्किक परिणाम भी है;
\item यदि !!{sentence}s का कोई समुच्चय असंगत है, तो वह असंतुष्टनीय है।
\end{enumerate}
ये !!a{derivation} तंत्र के महत्त्वपूर्ण गुण हैं। यदि इनमें से कोई गुण न हो,
तो !!{derivation} तंत्र दोषपूर्ण होगा---वह आवश्यकता से अधिक बातें !!{derive}
करेगा। इसलिए !!a{derivation} तंत्र की सुदृढ़ता स्थापित करना अत्यंत
महत्त्वपूर्ण है।
\end{explain}

\begin{thm}[सुदृढ़ता]
\ollabel{thm:soundness}
यदि $!A$ !!{derivable} है और इस व्युत्पत्ति की !!{undischarged} मान्यताएँ
$\Gamma$ हैं, तो $\Gamma \Entails !A$।
\end{thm}

\begin{proof}
$\delta$ को $!A$ की !!a{derivation} मानें। हम~$\delta$ में अनुमानों की
संख्या पर आगमन करते हैं।

आगमन-आधार के लिए, अनुमानों की संख्या~$0$ होने पर कथन सिद्ध करते हैं। इस
स्थिति में $\delta$ में केवल एक !!{sentence}~$!A$, अर्थात् एक मान्यता, होती
है। वह मान्यता !!{undischarged} है, क्योंकि मान्यताएँ केवल अनुमानों द्वारा ही
!!{discharged} हो सकती हैं और यहाँ कोई अनुमान नहीं है। अतः ऐसी प्रत्येक
\iftag{FOL}{!!{structure}~$\Struct{M}$}{!!{valuation}~$\pAssign{v}$}
जो प्रमाण की सभी !!{undischarged} मान्यताओं को संतुष्ट करती है,~$!A$ को भी
संतुष्ट करती है।

अब आगमन-चरण लें। मानें कि $\delta$ में~$n$ अनुमान हैं। सबसे निचले अनुमान के
पूर्वाधार !!{derive} किए गए हैं; इसके लिए उप-!!{derivation}s का उपयोग हुआ है,
और इनमें से प्रत्येक में~$n$ से कम अनुमान हैं। हम आगमन-परिकल्पना मानते हैं:
सबसे निचले अनुमान के पूर्वाधार !!{undischarged} मान्यताओं से तार्किकतः
निष्पन्न होते हैं; ये उन पूर्वाधारों पर समाप्त होने वाली उप-!!{derivation}s की
मान्यताएँ हैं। हमें दिखाना है कि
निष्कर्ष~$!A$ पूरे प्रमाण की !!{undischarged} मान्यताओं से तार्किकतः निष्पन्न
होता है।

हम सबसे निचले अनुमान के प्रकार के अनुसार प्रकरणों में भेद करते हैं। पहले केवल
एक पूर्वाधार वाले संभावित अनुमानों पर विचार करते हैं।

\begin{enumerate}
\item मानें कि अंतिम अनुमान \Intro{\lnot} है। !!{derivation} का रूप है
  \begin{prooftree}
    \AxiomC{$\Gamma, \Discharge{!A}{n}$}
    \RightLabel{$\delta_1$}
    \DeduceC{$\lfalse$}
    \DischargeRule{\Intro{\lnot}}{n}
    \UnaryInfC{$\lnot !A$}
  \end{prooftree}
  आगमन-परिकल्पना से $\lfalse$, !!{undischarged} मान्यताओं
  $\Gamma \cup \{!A\}$ से तार्किकतः निष्पन्न होता है; ये~$\delta_1$ की
  मान्यताएँ हैं। किसी
  \iftag{FOL}{!!a{structure}~$\Struct{M}$}{!!a{valuation}~$\pAssign{v}$} पर
  विचार करें। हमें दिखाना है कि यदि
  $\iftag{FOL}{\Sat{M}}{\pSat{v}}{\Gamma}$, तो
  $\iftag{FOL}{\Sat{M}}{\pSat{v}}{\lnot !A}$। विरोधाभासार्थ मानें कि
  $\iftag{FOL}{\Sat{M}}{\pSat{v}}{\Gamma}$, पर
  $\iftag{FOL}{\Sat/{M}}{\pSat/{v}}{\lnot !A}$, अर्थात्
  $\iftag{FOL}{\Sat{M}}{\pSat{v}}{!A}$। तब
  $\iftag{FOL}{\Sat{M}}{\pSat{v}}{\Gamma \cup \{!A\}}$ होगा। यह हमारी
  आगमन-परिकल्पना के विरुद्ध है। अतः
  $\iftag{FOL}{\Sat{M}}{\pSat{v}}{\lnot !A}$.

\item अंतिम अनुमान \Elim{\land} है। इसके दो रूप हैं: $!A$ या $!B$ को
  पूर्वाधार $!A \land !B$ से निष्कर्षित किया जा सकता है। पहले प्रकरण पर विचार करें।
  !!{derivation}~$\delta$ इस प्रकार है:
  \begin{prooftree}
    \AxiomC{$\Gamma$}
    \RightLabel{$\delta_1$}
    \DeduceC{$!A \land !B$}
    \RightLabel{\Elim{\land}}
    \UnaryInfC{$!A$}
  \end{prooftree}
  आगमन-परिकल्पना से $!A \land !B$, !!{undischarged} मान्यताओं~$\Gamma$ से
  तार्किकतः निष्पन्न होता है; ये~$\delta_1$ की मान्यताएँ हैं। किसी
  !!a{structure}~\iftag{FOL}{$\Struct{M}$}{$\pAssign{v}$} पर विचार करें। हमें
  दिखाना है कि यदि $\iftag{FOL}{\Sat{M}}{\pSat{v}}{\Gamma}$, तो
  $\iftag{FOL}{\Sat{M}}{\pSat{v}}{!A}$। मानें कि
  $\iftag{FOL}{\Sat{M}}{\pSat{v}}{\Gamma}$। हमारी आगमन-परिकल्पना
  ($\Gamma \Entails !A \land !B$) से हम जानते हैं कि
  $\iftag{FOL}{\Sat{M}}{\pSat{v}}{!A \land !B}$। परिभाषा से,
  $\iftag{FOL}{\Sat{M}}{\pSat{v}}{!A \land !B}$ तभी और केवल तभी जब
  $\iftag{FOL}{\Sat{M}}{\pSat{v}}{!A}$ और
  $\iftag{FOL}{\Sat{M}}{\pSat{v}}{!B}$। (जिस प्रकरण में $!B$ को
  $!A \land !B$ से निष्कर्षित किया जाता है, वह इसी प्रकार सँभलता है।)
  
\item अंतिम अनुमान \Intro{\lor} है। इसके दो रूप हैं: $!A
  \lor !B$ को
  पूर्वाधार~$!A$ या पूर्वाधार~$!B$ से निष्कर्षित किया जा सकता है। पहले प्रकरण
  पर विचार करें। !!{derivation} का रूप है
  \begin{prooftree}
    \AxiomC{$\Gamma$}
    \RightLabel{$\delta_1$}
    \DeduceC{$!A$}
    \RightLabel{\Intro{\lor}}
    \UnaryInfC{$!A \lor !B$}
  \end{prooftree}
  आगमन-परिकल्पना से $!A$, !!{undischarged} मान्यताओं~$\Gamma$ से तार्किकतः
  निष्पन्न होता है; ये~$\delta_1$ की मान्यताएँ हैं। किसी
  \iftag{FOL}{!!a{structure}~$\Struct{M}$}{!!a{valuation}~$\pAssign{v}$} पर
  विचार करें। हमें दिखाना है कि यदि
  $\iftag{FOL}{\Sat{M}}{\pSat{v}}{\Gamma}$, तो
  $\iftag{FOL}{\Sat{M}}{\pSat{v}}{!A \lor !B}$। मानें कि
  $\iftag{FOL}{\Sat{M}}{\pSat{v}}{\Gamma}$; तब
  $\iftag{FOL}{\Sat{M}}{\pSat{v}}{!A}$, क्योंकि $\Gamma \Entails !A$
  (आगमन-परिकल्पना)। अतः
  $\iftag{FOL}{\Sat{M}}{\pSat{v}}{!A \lor !B}$ भी होना चाहिए। (जिस प्रकरण
  में $!A
  \lor !B$ को~$!B$ से निष्कर्षित किया जाता है, वह इसी प्रकार सँभलता है।)
  
\item अंतिम अनुमान \Intro{\lif} है। $!A \lif !B$ को मान्यता~$!A$ और
  निष्कर्ष~$!B$ वाले उपप्रमाण से निष्कर्षित किया जाता है, अर्थात्
  \begin{prooftree}
    \AxiomC{$\Gamma, \Discharge{!A}{n}$}
    \RightLabel{$\delta_1$}
    \DeduceC{$!B$}
    \DischargeRule{\Intro{\lif}}{n}
    \UnaryInfC{$!A \lif !B$}
  \end{prooftree}
  आगमन-परिकल्पना से $!B$,~$\delta_1$ की !!{undischarged} मान्यताओं से
  तार्किकतः निष्पन्न होता है, अर्थात् $\Gamma \cup \{!A\} \Entails
  !B$। किसी
  \iftag{FOL}{!!a{structure}~$\Struct{M}$}{!!a{valuation}~$\pAssign{v}$} पर
  विचार करें।~$\delta$ की !!{undischarged} मान्यताएँ केवल $\Gamma$ हैं,
  क्योंकि अंतिम अनुमान पर $!A$ मुक्त कर दिया गया है। अतः हमें दिखाना है कि
  $\Gamma \Entails !A \lif !B$। विरोधाभासार्थ मानें कि किसी
  \iftag{FOL}{!!{structure}~$\Struct{M}$}{!!{valuation}~$\pAssign{v}$} के लिए
  $\iftag{FOL}{\Sat{M}}{\pSat{v}}{\Gamma}$, पर
  $\iftag{FOL}{\Sat/{M}}{\pSat/{v}}{!A \lif !B}$। तब
  $\iftag{FOL}{\Sat{M}}{\pSat{v}}{!A}$ और
  $\iftag{FOL}{\Sat/{M}}{\pSat/{v}}{!B}$। किन्तु परिकल्पना से $!B$,
  $\Gamma \cup \{!A\}$ का तार्किक परिणाम है, अर्थात्
  $\iftag{FOL}{\Sat{M}}{\pSat{v}}{!B}$, जो विरोधाभास है। अतः
  $\Gamma \Entails !A \lif !B$.

\item अंतिम अनुमान \FalseInt है। यहाँ $\delta$ का अंत इस प्रकार होता है:
  \begin{prooftree}
    \AxiomC{$\Gamma$}
    \RightLabel{$\delta_1$}
    \DeduceC{$\lfalse$}
    \RightLabel{\FalseInt}
    \UnaryInfC{$!A$}
  \end{prooftree}
  आगमन-परिकल्पना से $\Gamma \Entails \lfalse$। हमें दिखाना है कि
  $\Gamma \Entails !A$। मानें कि ऐसा नहीं है; तब किसी
  ~$\iftag{FOL}{\Struct{M}}{\pAssign{v}}$ के लिए
    $\iftag{FOL}{\Sat{M}}{\pSat{v}}{\Gamma}$ और
    $\iftag{FOL}{\Sat/{M}}{\pSat/{v}}{!A}$। पर हमारे पास सदैव
    $\iftag{FOL}{\Sat/{M}}{\pSat/{v}}{\lfalse}$ है, इसलिए इसका अर्थ
    $\Gamma \Entails/ \lfalse$ होगा, जो आगमन-परिकल्पना के विरुद्ध है।

\item अंतिम अनुमान \FalseCl है: अभ्यास।

\iftag{FOL}{
\item अंतिम अनुमान \Intro{\lforall} है। तब $\delta$ का रूप है
  \begin{prooftree}
    \AxiomC{$\Gamma$}
    \RightLabel{$\delta_1$}
    \DeduceC{$!A(a)$}
    \RightLabel{\Intro{\lforall}}
    \UnaryInfC{$\lforall[x][!A(x)]$}
  \end{prooftree}
  आगमन-परिकल्पना से पूर्वाधार $!A(a)$, !!{undischarged} मान्यताओं $\Gamma$
  का तार्किक परिणाम है। ऐसी किसी संरचना~$\Struct{M}$ पर विचार करें कि
  $\Sat{M}{\Gamma}$। हमें दिखाना है कि
  $\Sat{M}{\lforall[x][!A(x)]}$। चूँकि $\lforall[x][!A(x)]$ !!a{sentence}
  है, इसका अर्थ है कि हमें प्रत्येक चर-नियतन~$s$ के लिए
  $\Sat{M}{!A(x)}[s]$ दिखाना है
  (\olref[syn][ass]{prop:sat-quant})। चूँकि $\Gamma$ में केवल वाक्य हैं,
  $\Sat{M}{!B}[s]$ हर $!B \in \Gamma$ के लिए,
  \olref[syn][sat]{defn:satisfaction} के अनुसार।~$\Struct{M'}$ को
  $\Struct{M}$ जैसा ही मानें, सिवाय इसके
  कि $\Assign{a}{M'} = s(x)$। चूँकि $a$,~$\Gamma$ में नहीं आता,
  \olref[syn][ext]{cor:extensionality-sent} से $\Sat{M'}{\Gamma}$। चूँकि
  $\Gamma \Entails
  !A(a)$, इसलिए $\Sat{M'}{!A(a)}$। चूँकि $!A(a)$
  !!a{sentence} है, \olref[syn][ass]{prop:sentence-sat-true} से
  $\Sat{M'}{!A(a)}[s]$। \olref[syn][ext]{prop:ext-formulas} से
  $\Sat{M'}{!A(x)}[s]$ तभी और केवल तभी जब $\Sat{M'}{!A(a)}$ (याद रखें कि
  $!A(a)$ केवल $\Subst{!A(x)}{a}{x}$ है)। अतः
  $\Sat{M'}{!A(x)}[s]$। चूँकि $a$,~$!A(x)$ में नहीं आता,
  \olref[syn][ext]{prop:extensionality} से $\Sat{M}{!A(x)}[s]$। पर $s$ एक
  स्वेच्छ चर-नियतन था, अतः $\Sat{M}{\lforall[x][!A(x)]}$।
  
\item अंतिम अनुमान \Intro{\lexists} है: अभ्यास।

\item अंतिम अनुमान \Elim{\forall} है: अभ्यास।
}{}
\end{enumerate}

अब अनेक पूर्वाधारों वाले संभावित अनुमानों पर विचार करें:
\Elim{\lor}, \Intro{\land}, \iftag{FOL}{\Elim{\lif}, और
  \Elim{\lexists}}{और \Elim{\lif}}।
\begin{enumerate}

\item अंतिम अनुमान \Intro{\land} है। $!A \land !B$ को पूर्वाधारों $!A$
  और $!B$ से निष्कर्षित किया जाता है और $\delta$ का रूप है
  \begin{prooftree}
    \AxiomC{$\Gamma_1$}
    \RightLabel{$\delta_1$}
    \DeduceC{$!A$}
    \AxiomC{$\Gamma_2$}
    \RightLabel{$\delta_2$}
    \DeduceC{$!B$}
    \RightLabel{\Intro{\land}}
    \BinaryInfC{$!A \land !B$}
  \end{prooftree}
  आगमन-परिकल्पना से $!A$, !!{undischarged} मान्यताओं~$\Gamma_1$ से
  तार्किकतः निष्पन्न होता है; ये~$\delta_1$ की मान्यताएँ हैं। इसी प्रकार
  $!B$, !!{undischarged} मान्यताओं~$\Gamma_2$ से तार्किकतः निष्पन्न होता है;
  ये~$\delta_2$ की मान्यताएँ हैं।~$\delta$ की !!{undischarged} मान्यताएँ
  $\Gamma_1 \cup
  \Gamma_2$ हैं,
  इसलिए हमें दिखाना है कि $\Gamma_1 \cup \Gamma_2 \Entails
  !A \land !B$।
  ऐसी किसी \iftag{FOL}{!!a{structure}~$\Struct{M}$}{!!a{valuation}~$\pAssign{v}$}
  पर विचार करें जिसके लिए
  $\iftag{FOL}{\Sat{M}}{\pSat{v}}{\Gamma_1 \cup \Gamma_2}$। चूँकि
  $\iftag{FOL}{\Sat{M}}{\pSat{v}}{\Gamma_1}$, इसलिए
  $\iftag{FOL}{\Sat{M}}{\pSat{v}}{!A}$ होना चाहिए, क्योंकि
  $\Gamma_1 \Entails !A$; और चूँकि
  $\iftag{FOL}{\Sat{M}}{\pSat{v}}{\Gamma_2}$, इसलिए
  $\iftag{FOL}{\Sat{M}}{\pSat{v}}{!B}$, क्योंकि $\Gamma_2 \Entails
  !B$।
  दोनों को मिलाकर,
  $\iftag{FOL}{\Sat{M}}{\pSat{v}}{!A \land !B}$।
  
\item अंतिम अनुमान \Elim{\lor} है: अभ्यास।

\item अंतिम अनुमान \Elim{\lif} है। $!B$ को पूर्वाधारों $!A \lif !B$
  और~$!A$ से निष्कर्षित किया जाता है। !!{derivation}~$\delta$ इस प्रकार है:
  \begin{prooftree}
    \AxiomC{$\Gamma_1$}
    \RightLabel{$\delta_1$}
    \DeduceC{$!A \lif !B$}
    \AxiomC{$\Gamma_2$}
    \RightLabel{$\delta_2$}
    \DeduceC{$!A$}
    \RightLabel{\Elim{\lif}}
    \BinaryInfC{$!B$}
  \end{prooftree}
  आगमन-परिकल्पना से $!A \lif !B$, !!{undischarged} मान्यताओं~$\Gamma_1$ से
  तार्किकतः निष्पन्न होता है; ये~$\delta_1$ की मान्यताएँ हैं। इसी प्रकार
  $!A$, !!{undischarged} मान्यताओं~$\Gamma_2$ से तार्किकतः निष्पन्न होता है;
  ये~$\delta_2$ की मान्यताएँ हैं। किसी
  \iftag{FOL}{!!a{structure}~$\Struct{M}$}{!!a{valuation}~$\pAssign{v}$} पर
  विचार करें। हमें दिखाना है कि यदि
  $\iftag{FOL}{\Sat{M}}{\pSat{v}}{\Gamma_1 \cup
    \Gamma_2}$, तो
  $\iftag{FOL}{\Sat{M}}{\pSat{v}}{!B}$। मानें कि
  $\iftag{FOL}{\Sat{M}}{\pSat{v}}{\Gamma_1 \cup \Gamma_2}$। चूँकि
  $\Gamma_1 \Entails !A \lif !B$, इसलिए
  $\iftag{FOL}{\Sat{M}}{\pSat{v}}{!A
    \lif !B}$। चूँकि
  $\Gamma_2 \Entails !A$, हमारे पास
  $\iftag{FOL}{\Sat{M}}{\pSat{v}}{!A}$ है। इसका अर्थ है कि
  $\iftag{FOL}{\Sat{M}}{\pSat{v}}{!B}$ (क्योंकि यदि
  $\iftag{FOL}{\Sat/{M}}{\pSat/{v}}{!B}$, तो
  $\iftag{FOL}{\Sat{M}}{\pSat{v}}{!A}$ होने से
  $\iftag{FOL}{\Sat/{M}}{\pSat/{v}}{!A \lif !B}$ होता, जो
  ~$\iftag{FOL}{\Sat{M}}{\pSat{v}}{!A \lif !B}$ के विरुद्ध है)।

\item अंतिम अनुमान \Elim{\lnot} है: अभ्यास।
  
\tagitem{FOL}{अंतिम अनुमान \Elim{\lexists} है: अभ्यास।}{}
\end{enumerate}
\end{proof}

\tagprob{FOL}
\begin{prob}
\olref[fol][ntd][sou]{thm:soundness} का प्रमाण पूरा कीजिए।
\end{prob}
\tagendprob

\tagprob{notFOL}
\begin{prob}
\olref[pl][ntd][sou]{thm:soundness} का प्रमाण पूरा कीजिए।
\end{prob}
\tagendprob

\begin{cor}
\ollabel{cor:weak-soundness}
यदि $\Proves !A$, तो $!A$ \iftag{FOL}{वैध}{सर्वसत्य} है।
\end{cor}

\begin{cor}
\ollabel{cor:consistency-soundness}
यदि $\Gamma$ संतुष्टनीय है, तो वह संगत है।
\end{cor}

\begin{proof}
हम प्रतिधन कथन सिद्ध करते हैं। मान लें कि $\Gamma$ संगत नहीं है। तब
$\Gamma \Proves \lfalse$, अर्थात् $\lfalse$ की ऐसी कोई !!a{derivation} है
जिसकी !!{undischarged} मान्यताएँ~$\Gamma$ में हैं।
\olref{thm:soundness} से ऐसी प्रत्येक
\iftag{FOL}{!!{structure}~$\Struct{M}$}{!!{valuation}~$\pAssign{v}$}
जो $\Gamma$ को संतुष्ट करती है, उसे $\lfalse$ को भी संतुष्ट करना होगा। चूँकि
$\iftag{FOL}{\Sat/{M}}{\pSat/{v}}{\lfalse}$ प्रत्येक
\iftag{FOL}{!!{structure}~$\Struct{M}$}{!!{valuation}~$\pAssign{v}$} के लिए
सत्य है, इसलिए कोई
\iftag{FOL}{$\Struct{M}$}{$\pAssign{v}$},~$\Gamma$ को संतुष्ट नहीं कर सकती;
अर्थात् $\Gamma$ संतुष्टनीय नहीं है।
\end{proof}

\end{document}
